\Article{231}
The collection of \textit{all} individual characters of a given form or class makes up the complete character of that form or class. For \textit{example}, the complete character of the form $(10,3,17)$, or of the whole class it represents, is $1,4$; $N7$; $N23$. Similarly, the complete character of the form $(7,1,-17)$ is $7,8$; $R3$; $N5$, since the individual character $3,4$ is left out in this case because it is already contained in the character $7,8$. --- On this basis we subdivide the whole order of properly primitive classes (positive ones when the determinant is negative) with a given determinant into several different \textit{genera}: classes with the same complete character belong to the same genus, and classes with different complete characters to different genera. We assign each genus the complete character of the classes it contains. Thus, for \textit{example}, for the determinant $-161$ there are sixteen
\SourceLeaf{239}
positive properly primitive classes. They are divided into four genera as follows:
\begin{center}
\begin{tabular}{c|l}
Character & \multicolumn{1}{c}{Representative forms of the classes} \\
$1,4;\ R7;\ R23$ & $(1,0,161)$, $(2,1,81)$, $(9,1,18)$, $(9,-1,18)$ \\
$1,4;\ N7;\ N23$ & $(5,2,33)$, $(5,-2,33)$, $(10,3,17)$, $(10,-3,17)$ \\
$3,4;\ R7;\ N23$ & $(7,0,23)$, $(11,2,15)$, $(11,-2,15)$, $(14,7,15)$ \\
$3,4;\ N7;\ R23$ & $(3,1,54)$, $(3,-1,54)$, $(6,1,27)$, $(6,-1,27)$
\end{tabular}
\end{center}

For the number of different complete characters possible in principle, keep the following in mind.

I. When the determinant $D$ is divisible by $8$, four different individual characters are possible with respect to $8$; the number $4$ gives no separate character (note to the preceding article). In addition, each odd prime divisor of $D$ gives two characters. So if there are $m$ of these divisors, there are altogether $2^{m+2}$ different complete characters (setting $m=0$ whenever $D$ is a power of two).

II. When the determinant $D$ is not divisible by $8$, but is divisible by $4$ and also by $m$ odd primes, there are altogether $2^{m+1}$ different complete characters.

III. When the determinant $D$ is even but not divisible by $4$, it is either $\equiv2\ (\mathrm{mod.}\ 8)$ or $\equiv6$. In the first case there are two individual characters with respect to $8$, namely $1\ \mathit{and}\ 7,8$ and $3\ \mathit{and}\ 5,8$; in the second case there are just as many. Thus, setting the number of odd prime divisors of $D$ $=m$, there are altogether $2^{m+1}$ different complete characters.

IV. When $D$ is odd, it is either $\equiv1$ or $\equiv3\ (\mathrm{mod.}\ 4)$. In the latter case there are two different characters with respect to $4$, whereas in the former case no such relation enters the complete character. Thus, with $m$ denoting the same thing as before, there are $2^m$ different complete characters in the former case and $2^{m+1}$ in the latter.

But we must take care to note that this by no means implies that as many genera actually
\SourceLeaf{240}
exist as there are different characters possible in principle. In our example, only half of these actually correspond to classes or genera. There are no positive classes with the characters $1,4;\ R7;\ N23$ or $1,4;\ N7;\ R23$; or $3,4;\ R7;\ R23$ or $3,4;\ N7;\ N23$. We will discuss this very important subject more fully below.

From now on we call the form $(1,0,-D)$, which is undoubtedly the simplest of all forms with determinant $D$, the \textit{principal form}. The whole class containing it is called the \textit{principal class}, and the whole genus containing the principal class, the \textit{principal genus}. We must therefore distinguish carefully between the principal form, a form in the principal class, and a form in the principal genus; likewise between the principal class and a class in the principal genus. We always use these names even if, for some determinant, there are no classes other than the principal class, or no genera other than the principal genus. This happens, for \textit{example}, quite often when $D$ is a positive prime of the form $4n+1$.

\Needspace{8\baselineskip}
\Article{232}
Although our discussion of the characters of forms was introduced primarily to subdivide the \textit{positive properly primitive} order, nothing prevents us from applying the same ideas to negative forms and classes, or to improperly primitive ones. The positive improperly primitive order, the negative properly primitive order, and the negative improperly primitive order can all be subdivided into genera by the same principle. Thus, for \textit{example}, after the properly primitive order of forms with determinant $145$ has been divided into the following two genera,
\begin{center}
\begin{tabular}{l|l}
$R5,\ R29$ & $(1,0,-145)$, $(5,0,-29)$ \\
$N5,\ N29$ & $(3,1,-48)$, $(3,-1,-48)$
\end{tabular}
\end{center}
the improperly primitive order can likewise be divided into two genera:
\begin{center}
\begin{tabular}{l|l}
$R5,\ R29$ & $(4,1,-36)$, $(4,-1,-36)$ \\
$N5,\ N29$ & $(2,1,-72)$, $(10,5,-12)$
\end{tabular}
\end{center}
Or, just as the positive classes of forms with determinant $-129$ are divided into four genera,
\SourceLeaf{241}
\begin{center}
\begin{tabular}{l|l}
$1,4;\ R3;\ R43$ & $(1,0,129)$, $(10,1,13)$, $(10,-1,13)$ \\
$1,4;\ N3;\ N43$ & $(2,1,65)$, $(5,1,26)$, $(5,-1.\,26)$ \\
$3,4;\ R3;\ N43$ & $(3,0,43)$, $(7,2,19)$, $(7,-2,19)$ \\
$3,4;\ N3;\ R43$ & $(6,3,23)$, $(11,5,14)$, $(11,-5,14)$
\end{tabular}
\end{center}
the negative classes also fall into four orders:
\begin{center}
\begin{tabular}{l|l}
$3,4;\ N3;\ N43$ & $(-1,0,-129)$, $(-10,1,-13)$, $(-10,-1,-13)$ \\
$3,4;\ R3;\ R43$ & $(-2,1,-65)$, $(-5,1,-26)$, $(-5,-1,-26)$ \\
$1,4;\ N3;\ R43$ & $(-3,0,-43)$, $(-7,2,-19)$, $(-7,-2,-19)$ \\
$1,4;\ R3;\ N43$ & $(-6,3,-23)$, $(-11,5,-14)$, $(-11,-5,-14)$
\end{tabular}
\end{center}
But since the system of negative classes is always so similar to that of positive classes, it will usually seem unnecessary to construct it separately. Below we will show how to reduce the improperly primitive order to the properly primitive one.

Finally, no new rules are needed for subdividing derived orders. Every derived order comes from a primitive order (with a smaller determinant), and its individual classes correspond naturally to the individual classes of the latter. So the subdivision of the derived order can clearly be obtained from that of the primitive order.
